%==========================================================
%  CHAPITRE 4 — RÉALISATION
%==========================================================
\chapter{Réalisation}
\label{chap:realisation}

%──────────────────────────────────────────────────────────
\section{Environnement de développement}
\label{sec:real:env}
%──────────────────────────────────────────────────────────

\begin{center}
\begin{tabular}{>{\sffamily\bfseries\color{SecGrey}}l l l}
  \toprule
  \textbf{Catégorie} & \textbf{Outil / Technologie} & \textbf{Version} \\
  \midrule
  OS de développement   & \todo{Ubuntu 22.04 / Windows 11} & — \\
  Langage back-end      & \todo{Java / Python / Node.js}   & \todo{17 / 3.11 / 18} \\
  Framework back-end    & \todo{Spring Boot / FastAPI}      & \todo{3.2 / 0.110} \\
  Langage front-end     & \todo{JavaScript / TypeScript}    & \todo{ES2022} \\
  Framework front-end   & \todo{React / Vue / Angular}      & \todo{18.2} \\
  Base de données       & \todo{PostgreSQL / MySQL}         & \todo{15.4} \\
  ORM                   & \todo{Hibernate / SQLAlchemy}     & \todo{6.x} \\
  Conteneurisation      & \todo{Docker + Compose}           & \todo{24.0} \\
  Gestion de version    & \todo{Git + GitHub / GitLab}      & — \\
  CI/CD                 & \todo{GitHub Actions / Jenkins}   & — \\
  IDE                   & \todo{IntelliJ IDEA / VS Code}    & \todo{2024.1} \\
  \bottomrule
\end{tabular}
\end{center}

%──────────────────────────────────────────────────────────
\section{Implémentation — Module \todo{A}}
\label{sec:real:moduleA}
%──────────────────────────────────────────────────────────

\subsection{Description et fonctionnalités}

\todo{Décrire les fonctionnalités développées dans ce module,
les fichiers / packages concernés et les choix d'implémentation.}

\subsection{Extrait de code significatif}

L'extrait suivant illustre \todo{ce que fait ce code — ex : la logique
de création d'une commande avec validation métier}.

\begin{lstlisting}[
  language=Java,
  caption={\todo{Titre du snippet — ex : Service de création de commande}},
  label={lst:create-order}
]
// [A COMPLETER — coller ici un extrait de code representatif]
@Service
@Transactional
public class OrderService {

    private final OrderRepository orderRepository;
    private final InventoryClient inventoryClient;

    public OrderService(OrderRepository repo, InventoryClient inv) {
        this.orderRepository = repo;
        this.inventoryClient = inv;
    }

    public OrderDTO createOrder(CreateOrderRequest request) {
        // Valider la disponibilite en stock
        boolean inStock = inventoryClient.checkStock(
            request.getProductId(), request.getQuantity()
        );
        if (!inStock) {
            throw new InsufficientStockException(
                "Produit " + request.getProductId() + " indisponible."
            );
        }
        // Persister la commande
        Order order = Order.builder()
            .product(request.getProductId())
            .quantity(request.getQuantity())
            .status(OrderStatus.PENDING)
            .createdAt(Instant.now())
            .build();
        return OrderMapper.toDTO(orderRepository.save(order));
    }
}
\end{lstlisting}

\todo{Expliquer les choix de conception : pattern Repository, injection
de dépendances, gestion des exceptions, etc.}

%──────────────────────────────────────────────────────────
\section{Implémentation — Module \todo{B}}
\label{sec:real:moduleB}
%──────────────────────────────────────────────────────────

\subsection{Description et fonctionnalités}

\todo{Décrire les fonctionnalités du module B.}

\subsection{Extrait de code — \todo{exemple Python/JavaScript}}

\begin{lstlisting}[
  language=Python,
  caption={\todo{Titre du snippet Python}},
  label={lst:python-snippet}
]
# [A COMPLETER]
from fastapi import APIRouter, HTTPException, Depends
from sqlalchemy.orm import Session
from . import schemas, crud, database

router = APIRouter(prefix="/api/v1/items", tags=["items"])

@router.get("/{item_id}", response_model=schemas.ItemResponse)
def get_item(
    item_id: int,
    db: Session = Depends(database.get_db)
):
    """Recupere un item par son identifiant."""
    db_item = crud.get_item(db, item_id=item_id)
    if db_item is None:
        raise HTTPException(status_code=404, detail="Item non trouve")
    return db_item
\end{lstlisting}

%──────────────────────────────────────────────────────────
\section{Interfaces utilisateur}
\label{sec:real:ui}
%──────────────────────────────────────────────────────────

Cette section présente les principales interfaces de l'application.

\subsection{Page d'authentification}

\begin{figure}[H]
  \centering
  \placeholder{Capture d'écran : Page de connexion — figures/ui\_login.png}
  % \includegraphics[width=0.7\textwidth]{figures/ui_login}
  \caption{Interface d'authentification}
  \label{fig:ui-login}
\end{figure}

\todo{Décrire les éléments de l'interface : champs, validation, gestion
des erreurs, redirection après authentification.}

\subsection{Tableau de bord principal}

\begin{figure}[H]
  \centering
  \placeholder{Capture d'écran : Dashboard — figures/ui\_dashboard.png}
  % \includegraphics[width=\textwidth]{figures/ui_dashboard}
  \caption{Tableau de bord principal}
  \label{fig:ui-dashboard}
\end{figure}

\todo{Décrire les KPIs affichés, les graphiques, les actions rapides
disponibles depuis le dashboard.}

\subsection{\todo{Page / Module 3}}

\begin{figure}[H]
  \centering
  \placeholder{Capture d'écran : \todo{Nom de la page} — figures/ui\_page3.png}
  % \includegraphics[width=\textwidth]{figures/ui_page3}
  \caption{\todo{Titre de la figure}}
  \label{fig:ui-page3}
\end{figure}

%──────────────────────────────────────────────────────────
\section{Configuration et déploiement}
\label{sec:real:deploiement}
%──────────────────────────────────────────────────────────

\subsection{Docker Compose}

\begin{lstlisting}[
  language={},
  caption={Extrait du fichier docker-compose.yml},
  label={lst:docker-compose}
]
# [A COMPLETER — adapter a votre projet]
version: "3.9"
services:
  backend:
    build: ./backend
    ports:
      - "8080:8080"
    environment:
      - SPRING_DATASOURCE_URL=jdbc:postgresql://db:5432/mydb
      - SPRING_DATASOURCE_USERNAME=${DB_USER}
      - SPRING_DATASOURCE_PASSWORD=${DB_PASS}
    depends_on:
      - db

  frontend:
    build: ./frontend
    ports:
      - "3000:80"
    depends_on:
      - backend

  db:
    image: postgres:15-alpine
    volumes:
      - pgdata:/var/lib/postgresql/data
    environment:
      POSTGRES_DB: mydb
      POSTGRES_USER: ${DB_USER}
      POSTGRES_PASSWORD: ${DB_PASS}

volumes:
  pgdata:
\end{lstlisting}

\subsection{Pipeline CI/CD}

\begin{figure}[H]
  \centering
  \begin{tikzpicture}[
    node distance=0.5cm and 1.2cm,
    step/.style={
      rectangle, rounded corners=4pt,
      fill=ChapBlue!12, draw=ChapBlue,
      font=\sffamily\scriptsize\bfseries, text=ChapBlue,
      minimum width=2.2cm, minimum height=0.8cm,
      align=center
    },
    arrow/.style={-{Stealth[length=4pt]}, draw=AccentBlue, thick}
  ]
    \node[step] (push)  {Git Push};
    \node[step, right=of push] (build)  {Build};
    \node[step, right=of build] (test)  {Tests};
    \node[step, right=of test]  (scan)  {Analyse\\statique};
    \node[step, right=of scan]  (image) {Docker\\Image};
    \node[step, right=of image] (deploy){Déploiement};

    \draw[arrow] (push)  -- (build);
    \draw[arrow] (build) -- (test);
    \draw[arrow] (test)  -- (scan);
    \draw[arrow] (scan)  -- (image);
    \draw[arrow] (image) -- (deploy);
  \end{tikzpicture}
  \caption{Pipeline CI/CD — \todo{GitHub Actions / Jenkins}}
  \label{fig:cicd}
\end{figure}

\todo{Décrire chaque étape du pipeline : outil utilisé, durée moyenne,
critères de succès/échec.}

\clearpage
